\documentclass[11pt]{article}
\usepackage[T1]{fontenc}
\usepackage{lmodern,amsmath,amssymb,amsthm,microtype}
\usepackage[margin=1.03in]{geometry}
\usepackage[hidelinks]{hyperref}
\newtheorem{thm}{Theorem}
\newcommand{\av}{\operatorname{av}}
\newcommand{\gen}{\operatorname{Gen}}
\newcommand{\Res}{\operatorname{Res}}
\newcommand{\one}{\mathbf{1}}
\newcommand{\pn}{\psi_q}
\setlength{\emergencystretch}{2em}
\title{Detecting Dirichlet zeros with a sparse weighted count of prime powers}
\author{}\date{}
\begin{document}
\maketitle\vspace{-3.8em}
\begin{abstract}
Given a modulus $q$, select one generator from each nontrivial cyclic subgroup of its unit group. Weight each selected residue class by the number of generators of its subgroup. The resulting prime-power count, centered by its exact principal-character component, has square-root error up to logarithms if and only if the nonprincipal Dirichlet $L$-functions modulo $q$ have no zeros to the right of the critical line. The normalized cyclotomic trace of a pole residue measures the total multiplicity of a common zero. We classify all rational weights with this detection property in the finite multiplicity model and determine their minimum support when the identity class is excluded.
\end{abstract}

\section{One count for a fixed modulus}
Fix $q\ge3$. Write $G=(\mathbb Z/q\mathbb Z)^\times$, $n=|G|$, and let $e$ denote its identity. For every nontrivial cyclic subgroup $H\le G$, choose a generator $u_H$. Define
\[
 w(u_H)=\varphi(|H|),\qquad w(a)=0\quad(a\notin\{u_H\}_H).
\]
Every element other than $e$ generates exactly one such subgroup, so $\sum_a w(a)=n-1$. Extend weights by zero to nonunit classes. Using the von Mangoldt function $\Lambda$, set
\begin{align*}
 P_w(x)&=\sum_{k\le x}\Lambda(k)w(k\bmod q),\qquad
 \pn(x)=\sum_{\substack{k\le x\\(k,q)=1}}\Lambda(k),\\
 D_w(x)&=P_w(x)-\frac{n-1}{n}\pn(x).
\end{align*}
\begin{thm}\label{analytic}
The following conditions are equivalent for the chosen modulus.
\begin{enumerate}
\item Every nonprincipal Dirichlet $L(s,\chi)$ modulo $q$ is zero-free in $\Re s>1/2$.
\item $D_w(x)=O_q(\sqrt{x}\log^2(2x))$.
\item $D_w(x)=O_{q,\varepsilon}(x^{1/2+\varepsilon})$ for every $\varepsilon>0$.
\end{enumerate}
Assuming ordinary RH, conditions 2 and 3 retain their equivalence with 1 if $\pn(x)$ in the centering term is replaced by $x$.
\end{thm}
This equivalence takes the zero-free half-plane formulation used in the Lean repository \cite{lean}. Its forward estimate is the familiar prime number theorem in progressions under GRH. The choice of weights matters for the reverse implication, where zeros of several characters may occur at precisely the same complex number.

\section{Why common zeros survive the weighting}
For a rational weight $v$ on $G$ and $\chi\in\widehat G$, write
\[
 C_\chi(v)=\sum_{a\in G}v(a)\overline{\chi(a)}.
\]
Let $E=\exp G$, $K=\mathbb Q(\zeta_E)$, and let $\av:K\to\mathbb Q$ denote the field trace divided by $[K:\mathbb Q]$. Character values and all $C_\chi(v)$ lie in $K$. For $d\mid E$, the units modulo $E$ map onto the units modulo $d$ by the Chinese remainder theorem. Hence, for $H=\langle u\rangle$ of order $d$,
\[
 \varphi(d)\,\av(\overline{\chi(u)})=
 \sum_{a\in\gen(H)}\overline{\chi(a)},
\]
where $\gen(H)$ denotes the generators of $H$. Indeed, powers $u^j$ as $j$ runs over the units modulo $E$ visit each generator equally often. Summing over $H$ and using character orthogonality gives
\begin{equation}\label{trace}
 \av C_\chi(w)=\sum_{a\ne e}\overline{\chi(a)}
 =\begin{cases}n-1&\chi=\one,\\-1&\chi\ne\one.\end{cases}
\end{equation}
The coefficients may be complex and may differ between characters. Their normalized traces are identical on all nonprincipal characters.

\begin{proof}[Proof of Theorem~\ref{analytic}]
Write $\psi(x,\chi)=\sum_{k\le x}\Lambda(k)\chi(k)$. Fourier inversion gives
\[
 D_w(x)=\frac{1}{n}\sum_{\chi\ne\one} C_\chi(w)\psi(x,\chi).
\]
Under condition 1 the classical explicit formula for each nonprincipal character bounds its sum by $O_q(\sqrt x\log^2(2x))$ \cite{ep}. The omitted Euler factors of an imprimitive character change the sum by at most $O_q(\log(2x))$, since each prime dividing $q$ contributes at most $\log x$. This proves 2, which implies 3.

For $\Re s>1$, partial summation and the Euler products give
\begin{equation}\label{mellin}
 F(s)=s\int_1^\infty D_w(x)x^{-s-1}\,dx
 =-\frac1n\sum_{\chi\ne\one} C_\chi(w)\frac{L'}L(s,\chi).
\end{equation}
Condition 3 makes the integral locally uniformly convergent throughout $\Re s>1/2$ (on each compact set choose a smaller $\varepsilon$ than its distance to the boundary). Thus the meromorphic function on the right is holomorphic there. Suppose $\rho$ in this half-plane is a zero of one or more nonprincipal $L$-functions, of respective multiplicities $m_\chi\ge0$. Put $m_\chi=0$ for the other nonprincipal characters. All these $L$-functions are holomorphic at $\rho$. By \eqref{trace} their contribution satisfies the exact identity
\begin{equation}\label{residue}
 \av\!\left(\Res_{s=\rho}F(s)\right)
 =-\frac1n\sum_{\chi\ne\one}m_\chi\,\av C_\chi(w)
 =\frac1n\sum_{\chi\ne\one}m_\chi>0.
\end{equation}
This contradicts holomorphy and proves 1. Under RH the classical explicit formula gives $\pn(x)=x+O_q(\sqrt{x}\log^2(2x))$. Removing the primes dividing $q$ from the unrestricted Chebyshev sum costs $O_q(\log(2x))$. The final assertion follows.
\end{proof}

For $q=7$, choose $6,2,3$ as generators of the respective subgroups of orders $2,3,6$. Writing $\psi_a(x)=\sum_{k\le x,\ k\equiv a\ (7)}\Lambda(k)$, one obtains
\[
 P_w(x)=\psi_6(x)+2\psi_2(x)+2\psi_3(x).
\]
Thus the estimate $P_w(x)=5x/6+O(\sqrt{x}\log^2(2x))$, with ordinary RH supplied, detects every nonprincipal character modulo $7$. One may replace $2$ by $4$ or $3$ by $5$ independently. For cyclic $G$ of order $n$, there is one selected class for each divisor $d>1$ of $n$, hence $\tau(n)-1$ classes in total. If $t$ generates $G$, one convenient selection is $t^{n/d}$ for order $d$.

\section{Which rational weights detect every multiplicity pattern?}
Say that a weight $v$ detects multiplicities if $\sum_{\chi\ne\one}m_\chi C_\chi(v)\ne0$ for every nonzero tuple of nonnegative integers $m_\chi$. This finite definition allows every abstract pattern, independently of which coincidences can arise among actual Dirichlet zeros.
\begin{thm}\label{rational}
A rational weight $v$ detects multiplicities if and only if the rational numbers $b_\chi=\av C_\chi(v)$ for $\chi\ne\one$ are all strictly positive or all strictly negative.
\end{thm}
\begin{proof}
Either strict sign proves sufficiency by taking the normalized trace. For necessity, partition the nonprincipal characters into orbits under $\mathrm{Gal}(K/\mathbb Q)$. Rationality of $v$ gives $C_{\sigma\chi}(v)=\sigma C_\chi(v)$. On an orbit $O$ the trace $b_O$ is constant and $\sum_{\chi\in O}C_\chi(v)=|O|b_O\in\mathbb Q$, by averaging over the Galois group. An orbit with $b_O=0$ supplies an undetected pattern by giving all its members multiplicity one. If two orbits give opposite signs, positive integer multiples of their rational orbit sums cancel. Giving their members the corresponding multiplicities again supplies a nonzero undetected pattern.
\end{proof}
Consequently any rational weight satisfying the sign criterion admits the same analytic equivalences as Theorem~\ref{analytic}, with center $C_{\one}(v)\pn(x)/n$ and constants allowed to depend on $v$. The exact multiplicity formula \eqref{residue} uses the stronger normalization in \eqref{trace}.

\begin{thm}\label{support}
A set $A\subseteq G\setminus\{e\}$ supports a rational detecting weight if and only if $A$ meets $\gen(H)$ for every nontrivial cyclic subgroup $H$. The least support size is the number of nontrivial cyclic subgroups of $G$.
\end{thm}
\begin{proof}
Here a weight is supported on $A$ if it vanishes outside $A$. For sufficiency, take the generators in $A$ and use \eqref{trace}. For necessity suppose $A$ misses $B=\gen(H)$ for some nontrivial $H$. For every character put
\[
 m_\chi=|B|-\sum_{a\in B}\chi(a).
\]
The sum is an algebraic integer fixed by all cyclotomic automorphisms because these permute $B$ by powering. It is a rational integer of magnitude at most $|B|$. Thus all $m_\chi$ are nonnegative integers and $m_{\one}=0$. Orthogonality yields, for every $g\in G$,
\[
 \sum_\chi m_\chi\overline{\chi(g)}
 = n\bigl(|B|\one_{g=e}-\one_{g\in B}\bigr).
\]
Its value at $e$ is positive, so the tuple is nonzero. At each $g\in A$ the right side vanishes. Multiplication by any weights on $A$ and summation show that the tuple is undetected, even allowing complex weights. The generator sets are disjoint, so counting them gives the minimum.
\end{proof}
The lower bound concerns proofs that must distinguish all nonnegative multiplicity tuples using residues and that exclude $e$ from the weighted count. Allowing $e$ gives the singleton weight $v=\one_{\{e\}}$ with $C_\chi(v)=1$ for every $\chi$. Establishing which abstract tuples can occur as orders of common zeros is a separate arithmetic question. For the version centered at $\pn$, the support count concerns $P_w$. The centering total itself uses every reduced class.

\paragraph{Relation to earlier work.}
The repository \cite{lean} uses a one-sided estimate for a smoothed, quadratically twisted prime remainder and a pole argument. The RH overview \cite{wiki} describes how zero locations control prime-counting error. Here the estimate is two-sided at a fixed modulus. The distinguishing ingredient is the trace formula \eqref{residue} with a converse and a sharp support bound for all multiplicity patterns. Galois orbits of generators and character orthogonality are classical. A targeted search did not locate the combined criterion of Theorems~\ref{analytic}--\ref{support}. Historical priority remains open.

\begin{thebibliography}{9}\small
\bibitem{lean} L. Sernau et al., \emph{GRH for Dirichlet $L$-functions, in Lean 4}, readable edition, commit \texttt{72fead102f66}. \href{https://github.com/lsernau/grh-lean/tree/72fead102f663d5c7a407964a13dc0cf9d977207/grh}{Source and build notes}. See \texttt{src/New/Reduction/OneSidedBound.lean}.
\bibitem{ep} A.-M. Ernvall-Hyt\"onen and N. Paloj\"arvi, \emph{Explicit bound for the number of primes in arithmetic progressions assuming the Generalized Riemann Hypothesis}, Math. Comp. \textbf{91} (2022). \href{https://doi.org/10.1090/mcom/3691}{doi:10.1090/mcom/3691}.
\bibitem{wiki} Wikipedia contributors, \emph{Riemann hypothesis}, sections on the explicit formula and the distribution of primes, retrieved 26 November 2026. \url{https://en.wikipedia.org/wiki/Riemann_hypothesis}.
\end{thebibliography}
\end{document}
